\chapter{Τεχνολογίες και Πρότυπα}
\label{ch:technologies}

Το παρόν κεφάλαιο αναλύει σε βάθος τις τεχνολογίες και τα πρότυπα που χρησιμοποιήθηκαν. 
Για κάθε τεχνολογία παρουσιάζεται ο σκοπός της, ο τρόπος λειτουργίας της και ο ρόλος που διαδραματίζει στο σύνολο
του συστήματος. Η κατανόηση αυτών των τεχνολογιών αποτελεί προϋπόθεση για την
ανάγνωση των επόμενων κεφαλαίων.

% ==================================================================
\section{Κρυπτογραφία Δημόσιου Κλειδιού}
\label{sec:public-key-crypto}
% ==================================================================

Στο μοντέλο κρυπτογραφίας δημόσιου κλειδιού (\textlatin{asymmetric cryptography})
κάθε οντότητα κατέχει ένα ζεύγος κλειδιών:

\begin{itemize}
    \item Ιδιωτικό κλειδί (\textlatin{private key}): Παραμένει
    αυστηρά μυστικό και χρησιμοποιείται για τη δημιουργία ψηφιακών υπογραφών
    και την αποκρυπτογράφηση μηνυμάτων.
    \item Δημόσιο κλειδί (\textlatin{public key}): Κοινοποιείται
    ελεύθερα και χρησιμοποιείται για την επαλήθευση ψηφιακών υπογραφών και
    την κρυπτογράφηση μηνυμάτων.
\end{itemize}

Η σχέση μεταξύ των δύο κλειδιών είναι μαθηματικά αποδεδειγμένη: ό,τι υπογράφεται
με το ιδιωτικό κλειδί μπορεί να επαληθευτεί μόνο με το αντίστοιχο δημόσιο
κλειδί, και αντίστροφα. Αυτή η ιδιότητα επιτρέπει σε δύο μέρη που δεν
γνωρίζονται μεταξύ τους να εδραιώσουν εμπιστοσύνη χωρίς να χρειάζεται να
μοιραστούν κάποιο μυστικό εκ των προτέρων.

Στην  υλοποίηση μας χρησιμοποιείται ο αλγόριθμος \textlatin{RSA} με
μήκος κλειδιού 2048 \textlatin{bits}. Ο αλγόριθμος αυτός βασίζεται στη
δυσκολία παραγοντοποίησης μεγάλων ακεραίων αριθμών και αποτελεί τον πιο
ευρέως χρησιμοποιούμενο αλγόριθμο δημόσιου κλειδιού σε εφαρμογές ψηφιακών
υπογραφών.

% ==================================================================
\section{Ψηφιακά Πιστοποιητικά \textlatin{X.509 }  }
\label{sec:x509}
% ==================================================================

Η κρυπτογραφία δημόσιου κλειδιού λύνει το πρόβλημα της υπογραφής, αλλά
δημιουργεί ένα νέο, το πώς δηλαδή μπορεί κάποιος να βεβαιωθεί ότι ένα δημόσιο κλειδί
ανήκει σε συγκεκριμένη οντότητα και δεν έχει πλαστογραφηθεί. Η απάντηση
είναι τα ψηφιακά πιστοποιητικά \textlatin{X.509}~\cite{x509rfc} και η
Υποδομή Δημόσιου Κλειδιού (\textlatin{Public Key Infrastructure, PKI}).

Ένα πιστοποιητικό \textlatin{X.509} είναι ένα ψηφιακό έγγραφο που συνδέει
ένα δημόσιο κλειδί με μια ταυτότητα. Μπορεί να θεωρηθεί ως ένα ψηφιακό
διαβατήριο γιατί όπως ένα φυσικό διαβατήριο εκδίδεται από μια αρχή (το κράτος)
και πιστοποιεί την ταυτότητα του κατόχου, έτσι και ένα πιστοποιητικό
\textlatin{X.509} εκδίδεται από μια Αρχή Πιστοποίησης (\textlatin{Certificate
Authority, CA}) και πιστοποιεί ότι ένα δημόσιο κλειδί ανήκει σε συγκεκριμένη
οντότητα.

\subsection{Δομή Πιστοποιητικού}
\label{subsec:x509-structure}

Ένα πιστοποιητικό \textlatin{X.509 v3} αποτελείται από τα ακόλουθα πεδία:

\begin{itemize}
    \item Έκδοση (\textlatin{Version}): Η έκδοση του προτύπου
    (συνήθως \textlatin{v3}).
    \item Σειριακός Αριθμός (\textlatin{Serial Number}): Ένας
    μοναδικός αριθμός που εκχωρεί η \textlatin{CA} σε κάθε πιστοποιητικό
    που εκδίδει.
    \item Αλγόριθμος Υπογραφής (\textlatin{Signature Algorithm}):
    Ο αλγόριθμος που χρησιμοποίησε η \textlatin{CA} για να υπογράψει
    το πιστοποιητικό (π.χ. \textlatin{SHA-256 with RSA}).
    \item Εκδότης (\textlatin{Issuer}): Το \textlatin{Distinguished
    Name (DN)} της \textlatin{CA} που εξέδωσε το πιστοποιητικό.
    \item Περίοδος Ισχύος (\textlatin{Validity}): Τα πεδία
    \textlatin{notBefore} και \textlatin{notAfter} που ορίζουν το χρονικό
    παράθυρο εγκυρότητας.
    \item Υποκείμενο (\textlatin{Subject}): Το \textlatin{DN} του
    κατόχου. Στο πλαίσιο \textlatin{iSHARE}, το πεδίο \textlatin{Common
    Name (CN)} περιέχει ένα περιγραφικό όνομα, ενώ το πεδίο
    \textlatin{Serial Number} περιέχει το αναγνωριστικό \textlatin{EORI}
    του οργανισμού (π.χ. \textlatin{EU.EORI.NLPLEGMA}).
    \item Δημόσιο Κλειδί (\textlatin{Subject Public Key Info}):
    Ο αλγόριθμος και η τιμή του δημόσιου κλειδιού.
    \item Επεκτάσεις \textlatin{v3} (\textlatin{Extensions}):
    Πρόσθετες πληροφορίες, όπως:
    \begin{itemize}
        \item \textlatin{basicConstraints}: Υποδεικνύει αν το πιστοποιητικό
        ανήκει σε \textlatin{CA} (\textlatin{CA:TRUE}) ή σε τελικό χρήστη
        (\textlatin{CA:FALSE}).
        \item \textlatin{keyUsage}: Καθορίζει τις επιτρεπόμενες χρήσεις του
        κλειδιού. Η τιμή \textlatin{digitalSignature} είναι κρίσιμη για την
        υλοποίηση, καθώς ο \textlatin{Keyrock} ελέγχει ρητά αυτή την
        επέκταση κατά την επαλήθευση πιστοποιητικών. Αν απουσιάζει, ο
        \textlatin{Keyrock} απορρίπτει το πιστοποιητικό.
        \item \textlatin{subjectKeyIdentifier} και
        \textlatin{authorityKeyIdentifier}: Χρησιμοποιούνται για τη σύνδεση
        του πιστοποιητικού με τη \textlatin{CA} που το εξέδωσε.
    \end{itemize}
    \item Υπογραφή (\textlatin{Signature}): Η ψηφιακή υπογραφή
    της \textlatin{CA} πάνω σε όλα τα παραπάνω πεδία.
\end{itemize}

\subsection{Αλυσίδα Εμπιστοσύνης (\textlatin{Chain of Trust})}
\label{subsec:trust-chain}

Η εμπιστοσύνη σε ένα πιστοποιητικό μεταφράζεται σε εμπιστοσύνη στη
\textlatin{CA} που το εξέδωσε. Η \textlatin{CA} μπορεί με τη σειρά της
να είναι πιστοποιημένη από μια ανώτερη \textlatin{CA}, δημιουργώντας μια
ιεραρχία που ονομάζεται αλυσίδα εμπιστοσύνης (\textlatin{chain of trust}).

Η επαλήθευση μιας αλυσίδας εμπιστοσύνης γίνεται ως εξής:
\begin{enumerate}
    \item Ο παραλήπτης του πιστοποιητικού λαμβάνει το πιστοποιητικό μαζί με τα ενδιάμεσα πιστοποιητικά.
    \item Για κάθε πιστοποιητικό στην αλυσίδα, ελέγχει ότι η υπογραφή
    του έχει δημιουργηθεί από το ιδιωτικό κλειδί του αμέσως ανώτερου
    πιστοποιητικού.
    \item Η αλυσίδα τερματίζεται σε ένα πιστοποιητικό ρίζας
    (\textlatin{root certificate}) που ο επαληθευτής εμπιστεύεται εκ
    των προτέρων.
    \item Ελέγχεται ότι κανένα πιστοποιητικό στην αλυσίδα δεν έχει λήξει.
\end{enumerate}

Ως ψηφιακό αποτύπωμα (\textlatin{fingerprint}) ενός πιστοποιητικού ορίζεται 
η μοναδική συμβολοσειρά σταθερού μήκους που προκύπτει από την εφαρμογή μιας 
κρυπτογραφικής συνάρτησης κατακερματισμού, 
όπως ο αλγόριθμος \textlatin{SHA-256}, επί των δυαδικών δεδομένων του ίδιου 
του πιστοποιητικού. Το αποτύπωμα λειτουργεί ως ένα μοναδικό 
αναγνωριστικό το οποίο εξασφαλίζει την ακεραιότητα του πιστοποιητικού, καθώς 
οποιαδήποτε αλλαγή στα αρχικά δεδομένα θα παρήγαγε ένα εντελώς διαφορετικό 
αποτέλεσμα. Η χρήση του αποτυπώματος, αντί για την αποθήκευση ολόκληρου του αρχείου του πιστοποιητικού, 
βελτιστοποιεί τους πόρους του συστήματος και επιταχύνει τη διαδικασία 
επαλήθευσης της αλυσίδας εμπιστοσύνης.

Στην παρούσα υλοποίηση, η αλυσίδα αποτελείται από δύο μόνο επίπεδα:
\begin{itemize}
    \item Ρίζα (\textlatin{root}): Το πιστοποιητικό της τοπικής
    \textlatin{CA} (\textlatin{PlegmaCA}), του οποίου το αποτύπωμα
    (\textlatin{SHA-256 fingerprint}) καταχωρείται στη λίστα εμπιστευμένων
    \textlatin{CA} του \textlatin{iSHARE Satellite}.
    \item Φύλλα (\textlatin{leaf}): Τα πιστοποιητικά του παρόχου
    (\textlatin{EU.EORI.NLPLEGMA}) και του καταναλωτή
    (\textlatin{EU.EORI.NLPLEGMACONSUMER}), υπογεγραμμένα από τη
    \textlatin{PlegmaCA}.
\end{itemize}

\subsection{\textlatin{PKCS\#8} και \textlatin{PKCS\#1}}
\label{subsec:key-formats}

Ένα ζήτημα που αναδεικνύεται στην πράξη είναι η ύπαρξη δύο διαφορετικών
μορφών αποθήκευσης ιδιωτικών κλειδιών \textlatin{RSA}:

\begin{itemize}
    \item \textlatin{PKCS\#8}: Μια γενική μορφή που μπορεί να περιέχει κλειδιά
    οποιουδήποτε αλγορίθμου. Χρησιμοποιείται από το \textlatin{Keyrock} και το
    \textlatin{DSBA-PDP}. Η κεφαλίδα του είναι της μορφής \textlatin{-----BEGIN PRIVATE
    KEY-----}.
    \item \textlatin{PKCS\#1}: Μια ειδική μορφή αποκλειστικά για κλειδιά
    \textlatin{RSA}. Απαιτείται από τη βιβλιοθήκη
    \textlatin{lua-resty-jwt} που χρησιμοποιεί εσωτερικά το \textlatin{Kong
    API Gateway}. Η κεφαλίδα του είναι της μορφής \textlatin{-----BEGIN RSA
    PRIVATE KEY-----}.
\end{itemize}

Η διαφορά αυτή δεν είναι εμφανής θεωρητικά, αλλά στην πράξη οδηγεί σε
πολλά σφάλματα και στην υλοποιησή μας θα αναφέρουμε κάθε φορά τι χρησιμοποιείται
προκειμένου να αποφευχθούν.

% ==================================================================
\section{\textlatin{JSON Web Tokens (JWT)}}
\label{sec:jwt}
% ==================================================================

Τα \textlatin{JSON Web Tokens (JWT)}~\cite{rfc7519} αποτελούν το κύριο
μέσο ταυτοποίησης και εξουσιοδότησης στην αρχιτεκτονική
\textlatin{iSHARE}. Ένα \textlatin{JWT} είναι ένα αυτόνομο ψηφιακό
\textlatin{token} που μεταφέρει πληροφορίες, που ονομάζονται ισχυρισμοί, σε μορφή
\textlatin{JSON}, υπογεγραμμένο ψηφιακά ώστε ο παραλήπτης να μπορεί να
επαληθεύσει τόσο την ακεραιότητα όσο και τον εκδότη του.

\subsection{Δομή ενός \textlatin{JWT}}
\label{subsec:jwt-structure}

Ένα \textlatin{JWT} αποτελείται από τρία τμήματα, <\textlatin{Header}>.<\textlatin{Payload}>.<\textlatin{Signature}>, 
διαχωρισμένα με τελεία. Κάθε τμήμα κωδικοποιείται σε \textlatin{Base64URL}, μια παραλλαγή της
κωδικοποίησης \textlatin{Base64} ασφαλή για χρήση σε \textlatin{URL}.

Η επικεφαλίδα (\textlatin{Header}) δηλώνει τον τύπο του \textlatin{token} και τον αλγόριθμο
υπογραφής. Στο πλαίσιο \textlatin{iSHARE}, περιέχει επίσης την αλυσίδα
πιστοποιητικών \textlatin{X.509}:

\selectlanguage{english}
\begin{lstlisting}[language={}, caption={Header ενός iSHARE JWT.}]
{
  "alg": "RS256",
  "typ": "JWT",
  "x5c": [
    "<leaf_certificate_base64>",
    "<CA_certificate_base64>"
  ]
}
\end{lstlisting}
\selectlanguage{greek}

Τα πεδία ερμηνεύονται ως εξής:
\begin{itemize}
    \item \textlatin{alg}: Ο αλγόριθμος υπογραφής. Η τιμή \textlatin{RS256}
    σημαίνει υπογραφή \textlatin{RSA} με κατακερματισμό \textlatin{SHA-256}.
    Η διαδικασία περιλαμβάνει τον υπολογισμό του \textlatin{SHA-256 hash}
    των δύο πρώτων τμημάτων και την κρυπτογράφηση αυτού του \textlatin{hash}
    με το ιδιωτικό κλειδί \textlatin{RSA} του εκδότη.
    \item \textlatin{typ}: Ο τύπος \textlatin{token} που είναι πάντα \textlatin{JWT}.
    \item \textlatin{x5c}: Ένας πίνακας με τα
    πιστοποιητικά \textlatin{X.509} σε κωδικοποίηση \textlatin{Base64}
    (\textlatin{DER}). Το πρώτο στοιχείο είναι το πιστοποιητικό του εκδότη
    (\textlatin{leaf}), το δεύτερο το πιστοποιητικό της \textlatin{CA}.
    Αυτό επιτρέπει στον παραλήπτη να επαληθεύσει ολόκληρη την αλυσίδα
    εμπιστοσύνης χωρίς να χρειαστεί εξωτερική αναζήτηση.
\end{itemize}

Το ωφέλιμο φορτίο (\textlatin{Payload}) περιέχει τους ισχυρισμούς (\textlatin{claims}), δηλαδή
τις πληροφορίες που μεταφέρει το \textlatin{token}. Στο πλαίσιο
\textlatin{iSHARE}, χρησιμοποιούνται τα ακόλουθα τυποποιημένα
\textlatin{claims}:

\selectlanguage{english}
\begin{lstlisting}[language={}, caption={Payload ενός iSHARE JWT.}]
{
  "iss": "EU.EORI.NLPLEGMA",
  "sub": "EU.EORI.NLPLEGMA",
  "aud": "EU.EORI.NLPLEGMA",
  "jti": "a1b2c3d4-e5f6-7890-abcd-ef1234567890",
  "iat": 1716000000,
  "exp": 1716000030
}
\end{lstlisting}
\selectlanguage{greek}

Κάθε πεδίο αναλύεται ως εξής:
\begin{itemize}
    \item \textlatin{iss} (\textlatin{issuer}): Το αναγνωριστικό
    \textlatin{EORI} του εκδότη. Υποδεικνύει ποιος δημιούργησε και
    υπέγραψε το \textlatin{token}.
    \item \textlatin{sub} (\textlatin{subject}): Το αναγνωριστικό της
    οντότητας για την οποία εκδόθηκε το \textlatin{token} όπου συνήθως ταυτίζεται συνήθως με τον εκδότη.
    \item \textlatin{aud} (\textlatin{audience}): Ο αποδέκτης για τον οποίο
    προορίζεται το \textlatin{token}. Ο αποδέκτης πρέπει να ελέγξει ότι
    αυτό το πεδίο αντιστοιχεί στο δικό του αναγνωριστικό, αλλιώς να
    απορρίψει το \textlatin{token}. 
    \item \textlatin{jti} (\textlatin{JWT ID}): Ένα μοναδικό αναγνωριστικό
    σε μορφή \textlatin{UUID v4}. Αποτρέπει την επαναχρησιμοποίηση
    του ίδιου \textlatin{token} (\textlatin{replay attack}).
    \item \textlatin{iat} (\textlatin{issued at}): Η χρονοσφραγίδα
    \textlatin{Unix} (δευτερόλεπτα από 1/1/1970) στην οποία εκδόθηκε
    το \textlatin{token}.
    \item \textlatin{exp} (\textlatin{expiration}): Η χρονοσφραγίδα μετά
    την οποία το \textlatin{token} δεν είναι πλέον έγκυρο. Στο πλαίσιο
    \textlatin{iSHARE}, η τιμή αυτή ορίζεται σε 30 δευτερόλεπτα, ένα εξαιρετικά μικρό παράθυρο που ελαχιστοποιεί
    τον κίνδυνο κατάχρησης σε περίπτωση υποκλοπής.
\end{itemize}

Η υπογραφή (\textlatin{Signature})παράγεται εφαρμόζοντας τον αλγόριθμο \textlatin{RS256} στο
κείμενο που προκύπτει από τη συνένωση των δύο πρώτων τμημάτων
(κωδικοποιημένων σε \textlatin{Base64URL}) με τελεία:

\textlatin{\textlatin{RSASHA256(base64url(header) + "." + base64url(payload), private\_key)}}

Ο παραλήπτης επαληθεύει την υπογραφή εξάγοντας το δημόσιο κλειδί από
το πιστοποιητικό στο πεδίο \textlatin{x5c} και εφαρμόζοντας τον αντίστροφο
υπολογισμό. Αν η υπογραφή δεν ταιριάζει, το \textlatin{token}
απορρίπτεται.

\subsection{Διαφορές μεταξύ Παρόχου και Καταναλωτή \textlatin{JWT}}
\label{subsec:jwt-differences}

Μια κρίσιμη διάκριση στην υλοποίηση αφορά τον τρόπο με τον οποίο
δημιουργούνται τα \textlatin{JWT} ανάλογα με τον ρόλο του εκδότη:

\begin{itemize}
    \item Πάροχος: Τα πεδία \textlatin{iss}, \textlatin{sub} και
    \textlatin{aud} έχουν όλα την ίδια τιμή (\textlatin{EU.EORI.NLPLEGMA}).
    Ο πάροχος εκδίδει \textlatin{token} για τον εαυτό του, προς τον
    εαυτό του. Υπογράφει με το \textlatin{client.key} και ενσωματώνει
    στο \textlatin{x5c} τα \textlatin{client.pem} + \textlatin{ca.pem}.
    \item Καταναλωτής: Το πεδίο \textlatin{aud} δείχνει στον
    πάροχο (\textlatin{EU.EORI.NLPLEGMA}), υποδηλώνοντας ότι αυτό το
    \textlatin{token} προορίζεται για τον \textlatin{Keyrock} του
    παρόχου. Υπογράφει με το \textlatin{consumer.key} και ενσωματώνει
    στο \textlatin{x5c} τα \textlatin{consumer.pem} + \textlatin{ca.pem}.
\end{itemize}

% ==================================================================
\section{Το Πλαίσιο \textlatin{iSHARE}}
\label{sec:ishare}
% ==================================================================

Το \textlatin{iSHARE}~\cite{ishare2023} είναι ένα πλαίσιο εμπιστοσύνης
(\textlatin{trust framework}) που ορίζει τους κανόνες, τα πρωτόκολλα
και τις τεχνικές προδιαγραφές για τον ασφαλή διαμοιρασμό δεδομένων
μεταξύ οργανισμών. Αναπτύχθηκε αρχικά στις Κάτω Χώρες για τον κλάδο
της εφοδιαστικής αλυσίδας και υιοθετήθηκε από τη \textlatin{FIWARE
Foundation} ως βάση του πλαισίου \textlatin{i4Trust}~\cite{i4trust2023}
για τη δημιουργία \textlatin{Data Spaces}.

\subsection{Αναγνωριστικό \textlatin{EORI}}
\label{subsec:eori}

Κάθε συμμετέχων σε ένα \textlatin{iSHARE Data Space} αναγνωρίζεται
μοναδικά μέσω του αναγνωριστικού \textlatin{Economic Operators Registration
and Identification (EORI)}. Το \textlatin{EORI} είναι ένα αναγνωριστικό
που χρησιμοποιείται στην ευρωπαϊκή τελωνειακή νομοθεσία για την
αναγνώριση νομικών οντοτήτων που δραστηριοποιούνται σε εξωτερικό εμπόριο.

Αυτό το μοναδικό αναγνωριστικό καταχωρείται τόσο στο πιστοποιητικό
\textlatin{X.509} (στο πεδίο \textlatin{Serial Number} του \textlatin{DN})
όσο και στο μητρώο του \textlatin{Satellite}. Στην παρούσα υλοποίηση:
\begin{itemize}
    \item Πάροχος: \textlatin{EU.EORI.NLPLEGMA}
    \item Καταναλωτής: \textlatin{EU.EORI.NLPLEGMACONSUMER}
\end{itemize}

\subsection{Ροή Επαλήθευσης \textlatin{iSHARE} }
\label{subsec:ishare-flow-detail}

Η ροή επαλήθευσης κατά τον \textlatin{iSHARE} είναι μια ακολουθία
βημάτων που εξασφαλίζει ότι ο αιτών είναι αυτός που ισχυρίζεται
ότι είναι, ότι είναι αξιόπιστος συμμετέχων του \textlatin{Data
Space} και ότι διαθέτει τα κατάλληλα δικαιώματα πρόσβασης.
Αναλυτικά:

\begin{enumerate}
    \item Ο αιτών δημιουργεί \textlatin{iSHARE JWT}: Φτιάχνει
    ένα \textlatin{JWT} με τα στοιχεία του (πεδία \textlatin{iss},
    \textlatin{sub}, \textlatin{aud}) και ενσωματώνει την αλυσίδα
    πιστοποιητικών στο πεδίο \textlatin{x5c}. Υπογράφει με το ιδιωτικό
    κλειδί του.
    \item Ο αποδέκτης λαμβάνει το \textlatin{JWT}: Εξάγει
    το πιστοποιητικό \textlatin{leaf} από το πρώτο στοιχείο του
    \textlatin{x5c} και ελέγχει ότι η υπογραφή του \textlatin{JWT}
    αντιστοιχεί στο δημόσιο κλειδί αυτού του πιστοποιητικού.
    \item Έλεγχος αλυσίδας πιστοποιητικών: Ο αποδέκτης
    ελέγχει ότι το πιστοποιητικό \textlatin{leaf} έχει υπογραφεί
    από τη \textlatin{CA} και ότι αυτή η \textlatin{CA} βρίσκεται στη λίστα εμπιστευμένων
    \textlatin{CA} που παρέχει ο \textlatin{Satellite}
    (\textlatin{trusted\_list}).
    \item Επικοινωνία με τον \textlatin{Satellite}: Ο
    αποδέκτης αποκτά πρώτα ένα δικό του \textlatin{access token}
    από τον \textlatin{Satellite} (\textlatin{POST /token}) και στη
    συνέχεια χρησιμοποιεί αυτό το \textlatin{token} για να ζητήσει
    τη λίστα εμπιστευμένων \textlatin{CA}
    (\textlatin{GET /trusted\_list}) και τα στοιχεία του αιτούντα
    (\textlatin{GET /parties}).
    \item Έλεγχος κατάστασης συμμετέχοντα: Ο
    \textlatin{Satellite} επιστρέφει τα στοιχεία του αιτούντα
    (κατάσταση, πιστοποιήσεις, αποτύπωμα πιστοποιητικού). Ο
    αποδέκτης ελέγχει ότι η κατάσταση είναι \textlatin{Active} και
    ότι το αποτύπωμα αντιστοιχεί στο πιστοποιητικό του \textlatin{JWT}.
    \item Έλεγχος εξουσιοδότησης: Αν ο αιτών δεν είναι ο
    ίδιος ο πάροχος, ο αποδέκτης αναζητά πολιτική εκχώρησης
    εξουσιοδότησης (\textlatin{delegation evidence}) στο
    \textlatin{Authorization Registry}.
\end{enumerate}

\subsection{Πολιτικές Εκχώρησης Εξουσιοδότησης (\textlatin{Delegation Policies})}
\label{subsec:delegation}

Ο μηχανισμός εκχώρησης εξουσιοδότησης (\textlatin{delegation}) επιτρέπει
στον πάροχο δεδομένων να ορίζει ρητά τι δικαιούται να κάνει κάθε
καταναλωτής. Μια πολιτική εκχώρησης αποτελείται από τα ακόλουθα στοιχεία:

\begin{itemize}
    \item \textlatin{policyIssuer}: Ο οργανισμός που εκδίδει
    την πολιτική, πάντα ο πάροχος δεδομένων.
    \item \textlatin{target.accessSubject}: Ο αποδέκτης της
    πολιτικής, ο καταναλωτής που λαμβάνει δικαίωμα πρόσβασης.
    Στην υλοποίηση, χρησιμοποιείται το εσωτερικό \textlatin{UUID}
    του καταναλωτή στη βάση δεδομένων του \textlatin{Keyrock}.
    \item \textlatin{notBefore} / \textlatin{notOnOrAfter}:
    Χρονικό παράθυρο εγκυρότητας, σε μορφή \textlatin{Unix timestamp}.
    \item \textlatin{policySets}: Ένα ή περισσότερα σύνολα
    πολιτικών, καθένα περιέχοντας:
    \begin{itemize}
        \item \textlatin{resource.type}: Ο τύπος πόρου στον οποίο
        αφορά η πολιτική.
        \item \textlatin{resource.identifiers}: Ποιες οντότητες
        αφορά (\textlatin{["*"]} σημαίνει όλες).
        \item \textlatin{resource.attributes}: Ποια χαρακτηριστικά
        αφορά (\textlatin{["*"]} σημαίνει όλα).
        \item \textlatin{actions}: Ποιες ενέργειες επιτρέπονται
        (\textlatin{["GET"]}, \textlatin{["POST"]}, κ.λπ.).
        \item \textlatin{rules}: Η απόφαση (\textlatin{Permit} ή
        \textlatin{Deny}).
    \end{itemize}
\end{itemize}

Αυτή η δομή επιτρέπει ακριβή καθορισμό δικαιωμάτων (\textlatin{fine grained access control}).

% ==================================================================
\section{\textlatin{OAuth 2.0 }}
\label{sec:oauth2} 
% ==================================================================

Το \textlatin{OAuth 2.0} είναι ένα πρωτόκολλο
εξουσιοδότησης που ορίζει τον τρόπο με τον οποίο μια εφαρμογή μπορεί
να αποκτήσει πρόσβαση σε πόρους εκ μέρους ενός χρήστη ή οργανισμού.
Δεν πρόκειται για πρωτόκολλο ταυτοποίησης αλλά ασχολείται αποκλειστικά
με την εξουσιοδότηση.

Το \textlatin{OAuth 2.0} ορίζει τέσσερις βασικές ροές (\textlatin{grant
types}):

\begin{itemize}
    \item \textlatin{Authorization Code}: Για εφαρμογές ιστού
    (\textlatin{web apps}) όπου ο χρήστης ταυτοποιείται μέσω
    \textlatin{browser}. Είναι η πιο ασφαλής ροή.
    \item \textlatin{Client Credentials}: Για επικοινωνία μεταξύ
    μηχανών (\textlatin{M2M}), χωρίς εμπλοκή ανθρώπινου χρήστη. Αυτή
    η ροή χρησιμοποιείται στην παρούσα υλοποίηση.
    \item \textlatin{Resource Owner Password Credentials}:
    Για απευθείας αποστολή ονόματος χρήστη και κωδικού.
    \item \textlatin{Implicit}: Για εφαρμογές \textlatin{browser}
    χωρίς \textlatin{backend}, ροή που έχει αρχίσει να αποσύρεται.
\end{itemize}

\subsection{Ροή \textlatin{Client Credentials }}
\label{subsec:client-credentials}

Στην παρούσα υλοποίηση, ο καταναλωτής χρησιμοποιεί τη ροή
\textlatin{Client Credentials} για να αποκτήσει \textlatin{access
token} από τον \textlatin{Keyrock}. Η ροή αυτή προσαρμόζεται ώστε
αντί για παραδοσιακά \textlatin{client\_id, client\_secret} να
χρησιμοποιεί \textlatin{iSHARE JWT}.

Η αίτηση αποστέλλεται ως \textlatin{HTTP POST}:

\selectlanguage{english}
\begin{lstlisting}[language={}, caption={Αίτηση OAuth 2.0 Client Credentials με iSHARE JWT.}]
POST /oauth2/token HTTP/1.1
Content-Type: application/x-www-form-urlencoded

grant_type=client_credentials
&scope=iSHARE
&client_assertion_type=urn:ietf:params:oauth:client-assertion-type:jwt-bearer
&client_assertion=eyJhbGciOiJSUzI1NiIs...
&client_id=EU.EORI.NLPLEGMACONSUMER
\end{lstlisting}
\selectlanguage{greek}

Τα ορίσματα αναλύονται ως εξής:
\begin{itemize}
    \item \textlatin{grant\_type}: Η ροή \textlatin{OAuth 2.0} που
    χρησιμοποιείται. Η τιμή \textlatin{client\_credentials} υποδεικνύει
    ότι δεν εμπλέκεται ανθρώπινος χρήστης.
    \item \textlatin{scope}: Η τιμή \textlatin{iSHARE} ενεργοποιεί τη
    λογική \textlatin{iSHARE} στον \textlatin{Keyrock}, αντί της
    τυπικής ροής \textlatin{OAuth 2.0}.
    \item \textlatin{client\_assertion\_type}: Δηλώνει ότι τα
    διαπιστευτήρια του πελάτη παρέχονται ως \textlatin{JWT}
    (\textlatin{jwt-bearer}) αντί για \textlatin{client\_secret}.
    \item \textlatin{client\_assertion}: Το ίδιο το \textlatin{iSHARE
    JWT} του καταναλωτή, κωδικοποιημένο σε \textlatin{Base64URL}.
    \item \textlatin{client\_id}: Το \textlatin{EORI} του καταναλωτή.
\end{itemize}

Εφόσον η ταυτοποίηση του καταναλωτή ολοκληρωθεί επιτυχώς και επιβεβαιωθεί 
ότι διαθέτει τα απαραίτητα δικαιώματα πρόσβασης, ο διαχειριστής ταυτοτήτων 
(όπως το \textlatin{Keyrock}) εκδίδει ένα ψηφιακό διαπιστευτήριο (\textlatin{access token}).
Το διαπιστευτήριο αυτό λειτουργεί ως ένα ασφαλές διαβατήριο, καθώς εμπεριέχει την 
κρυπτογραφική απόδειξη των δικαιωμάτων που έχουν εκχωρηθεί στον χρήστη 
(\textlatin{delegation evidence}). Στη συνέχεια, ο καταναλωτής επισυνάπτει αυτό 
το διαπιστευτήριο σε κάθε αίτημά του προς την κεντρική πύλη ελέγχου του συστήματος 
(\textlatin{API Gateway}, ρόλο τον οποίο αναλαμβάνει το \textlatin{Kong}), 
προκειμένου να αποδείξει με ασφάλεια την εξουσιοδότησή του και να του 
επιτραπεί η τελική προσπέλαση στα δεδομένα.

% ==================================================================
\section{\textlatin{NGSI-LD} και \textlatin{JSON-LD}}
\label{sec:ngsi-ld}
% ==================================================================

Το \textlatin{NGSI-LD}~\cite{ngsi-ld2023} (\textlatin{Next Generation
Service Interfaces - Linked Data}) είναι ένα πρότυπο \textlatin{API}
που αναπτύχθηκε από τον ευρωπαϊκό οργανισμό τυποποίησης
\textlatin{ETSI} (προδιαγραφή \textlatin{ETSI GS CIM 009}). Αποτελεί
εξέλιξη του \textlatin{NGSI-v2} με ενσωμάτωση αρχών
\textlatin{Linked Data} μέσω \textlatin{JSON-LD}~\cite{jsonld2020}.

Το μοντέλο πληροφορίας του \textlatin{NGSI-LD} βασίζεται σε τρεις
θεμελιώδεις έννοιες:

\begin{itemize}
    \item Οντότητα (\textlatin{Entity}): Αναπαριστά ένα
    πραγματικό ή εννοιολογικό αντικείμενο (π.χ. ένα κτίριο, μια
    μέτρηση, μια συσκευή). Κάθε οντότητα έχει ένα μοναδικό
    αναγνωριστικό σε μορφή \textlatin{URI} και έναν τύπο.
    \item Ιδιότητα (\textlatin{Property}): Αναπαριστά ένα
    μετρήσιμο ή περιγραφικό χαρακτηριστικό μιας οντότητας. Κάθε
    ιδιότητα έχει μια τιμή (\textlatin{value}) και προαιρετικά μια
    μονάδα μέτρησης (\textlatin{unitCode}) και μια χρονοσφραγίδα
    (\textlatin{observedAt}).
    \item Σχέση (\textlatin{Relationship}): Αναπαριστά μια
    σύνδεση μεταξύ δύο οντοτήτων. Σε αντίθεση με τις ιδιότητες, η
    τιμή μιας σχέσης δεν είναι δεδομένο αλλά ένα \textlatin{URI}
    που δείχνει σε άλλη οντότητα.
\end{itemize}

\subsection{\textlatin{JSON-LD Context}}
\label{subsec:json-ld-context}

Το \textlatin{JSON-LD} (\textlatin{JSON for Linked Data})~\cite{jsonld2020}
είναι μια μέθοδος κωδικοποίησης \textlatin{Linked Data} χρησιμοποιώντας
τη μορφή \textlatin{JSON}. Η βασική ιδέα είναι η χρήση ενός ειδικού
πεδίου \textlatin{@context} που μετατρέπει τα σύντομα ονόματα ιδιοτήτων
σε παγκόσμια \textlatin{URI}.

Η χρησιμότητα αυτού γίνεται κατανοητή με ένα παράδειγμα: η ιδιότητα
\textlatin{activePower} σε ένα \textlatin{JSON} δεν φέρει εγγενώς σημασία,
κάθε οργανισμός θα μπορούσε να χρησιμοποιεί αυτό το όνομα με
διαφορετική σημασία. Ωστόσο, μέσω του \textlatin{@context}, η ιδιότητα
\textlatin{activePower} αντιστοιχίζεται στο \textlatin{URI}
\textlatin{https://plegma.example.org/vocab\#activePower}, το οποίο
είναι παγκόσμια μοναδικό.

\subsection{Βασικές Λειτουργίες \textlatin{API}}
\label{subsec:ngsi-api}

Το \textlatin{NGSI-LD API} ακολουθεί τις αρχές \textlatin{REST} και
παρέχει τις ακόλουθες βασικές λειτουργίες:

\begin{itemize}
    \item \textlatin{POST /ngsi-ld/v1/entities}: Δημιουργία νέας
    οντότητας.
    \item \textlatin{GET /ngsi-ld/v1/entities}: Αναζήτηση οντοτήτων με
    φίλτρα.
    \item \textlatin{GET /ngsi-ld/v1/entities/\{id\}}: Ανάκτηση
    συγκεκριμένης οντότητας μέσω αναγνωριστικού.
    \item \textlatin{PATCH /ngsi-ld/v1/entities/\{id\}/attrs}:
    Ενημέρωση ιδιοτήτων υπάρχουσας οντότητας.
    \item \textlatin{DELETE /ngsi-ld/v1/entities/\{id\}}: Διαγραφή
    οντότητας.
\end{itemize}

\subsection{\textlatin{Smart Data Models}}
\label{subsec:smart-data-models}

Τα \textlatin{Smart Data Models}~\cite{smart-data-models2023}
αποτελούν μια πρωτοβουλία που παρέχει ανοιχτά, τυποποιημένα σχήματα
δεδομένων για τα πιο κοινά αντικείμενα σε τομείς όπως η ενέργεια, τα
κτίρια, ο καιρός, οι μεταφορές κ.ά. Η πρωτοβουλία διοικείται
από τη \textlatin{FIWARE Foundation}, τον \textlatin{IUDX}, το
\textlatin{TM Forum} και τον \textlatin{OASC}.

Στην παρούσα εργασία, τα μοντέλα δεδομένων του \textlatin{Plegma
Dataset} σχεδιάστηκαν ώστε να κληρονομούν τη σημασιολογία των
επίσημων \textlatin{Smart Data Models}, επεκτείνοντάς τα με ιδιότητες
ειδικές για τα δεδομένα νοικοκυριών.

% ==================================================================
\section{Αρχιτεκτονική Ελέγχου Πρόσβασης \textlatin{XACML}}
\label{sec:xacml}
% ==================================================================

Η \textlatin{eXtensible Access Control Markup Language
(XACML)}~\cite{xacml3} είναι ένα πρότυπο \textlatin{OASIS} που ορίζει
τόσο μια γλώσσα έκφρασης πολιτικών ελέγχου πρόσβασης όσο και μια
αρχιτεκτονική αξιολόγησής τους. Η αρχιτεκτονική αυτή βασίζεται στον
διαχωρισμό ευθυνών μεταξύ τεσσάρων σημείων:

\begin{itemize}
    \item \textlatin{Policy Enforcement Point (PEP)}: Το
    σημείο που παρεμβάλλεται μεταξύ αιτούντα και πόρου. Λαμβάνει
    κάθε αίτηση, ρωτά τον \textlatin{PDP} αν πρέπει να επιτραπεί
    και εφαρμόζει την απόφαση. Στην παρούσα υλοποίηση, ο ρόλος αυτός
    ανατίθεται στο \textlatin{Kong API Gateway}.
    \item \textlatin{Policy Decision Point (PDP)}: Το
    σημείο που αξιολογεί μια αίτηση πρόσβασης σε σχέση με τις
    ισχύουσες πολιτικές. Λαμβάνει ως είσοδο τα στοιχεία της
    αίτησης (ποιος ζητά τι, πώς) και τις πολιτικές, και επιστρέφει
    απόφαση \textlatin{Permit} ή \textlatin{Deny}. Τον ρόλο αυτό
    επιτελεί ο \textlatin{DSBA-PDP} σε συνεργασία με τη λογική
    του \textlatin{plugin} \textlatin{ngsi-ishare-policies}.
    \item \textlatin{Policy Administration Point (PAP)}:
    Το σημείο όπου δημιουργούνται και αποθηκεύονται οι πολιτικές.
    Στην υλοποίηση, ο ρόλος αυτός ανατίθεται στον \textlatin{Keyrock}
    μέσω του \textlatin{endpoint} \textlatin{POST /ar/policy}.
    \item \textlatin{Policy Information Point (PIP)}:
    Το σημείο που παρέχει πρόσθετα δεδομένα που χρειάζεται ο
    \textlatin{PDP} για τη λήψη απόφασης. Στην υλοποίηση, αυτός
    ο ρόλος εκπληρώνεται μερικώς από τον \textlatin{iSHARE
    Satellite} (παρέχει πληροφορίες εμπιστοσύνης) και τον
    \textlatin{Keyrock} (παρέχει πολιτικές εκχώρησης).
\end{itemize}

Ο διαχωρισμός αυτός είναι σημαντικός γιατί επιτρέπει τη χρήση
διαφορετικών υλοποιήσεων για κάθε σημείο, εφόσον τηρούν τις
τυποποιημένες διεπαφές. Στην παρούσα αρχιτεκτονική, η λογική αξιολόγησης
πολιτικών \textlatin{iSHARE} (\textlatin{delegation evidence}) ενσωματώνεται
στο \textlatin{Kong plugin}, ενώ ο εξωτερικός \textlatin{DSBA-PDP}
παρέχει πρόσθετες δυνατότητες για πιο σύνθετα σενάρια.

